% !TeX root = ../papers/nc-p3.tex
\section{六次残余の最高重複度と各三重印での極曲線}
\label{sec:sextic-residual-triple-polars}
\begin{theorem}\label{thm:sextic-residual-triple-polars}
四点が共線でない相異なる十三点以上を三重以上に通る整六次曲面 $G$ は、どの点でも重複度三以下である。
各三重印 $p$ について、$G$ と $p$ に関する一階・二階の極曲面の共通零集合は曲線を含む。
特に前節の重複平面の六次残余の全ての平面外印に両結論が適用される。
\end{theorem}
\begin{proof}[次数二の射影の共通方向]
重複度四以上の点 $p$ の有無を調べる。六重点なら平面六次上の錐であり、
底の三重点は算術種数十と一つにつき少なくとも三の種数損失から三点以下である。
対応する三母線と頂点には十印までしか載らない。
五重点なら $p$ からの方向 $z$ で式は $tA_5(z)+A_6(z)$、両式は互いに素である。
一本の方向に $p$ 以外の三重印が $k$ 個あれば $1\le k\le3$。
一つなら $t$ をその点に合わせてずらして $A_5$ の位数は二以上、定数項は三以上となり交数は六以上。
二つ以上なら次数二以下の方向に関する局所Taylor係数のうち二次係数まで全て零である。
実際各係数は $t$ の一次以下で、二次係数でも両印で零になる必要がある。
ゆえに交数は九以上で、どの場合も $3k$ 以上である。
全交数三十より $p$ 以外に高々十印となる。

四重点なら式を $t^2 A_4(z)+t A_5(z)+A_6(z)$ とする。添字は方向平面での次数である。
$B=\gcd(A_4,A_5)$、$b=\deg B$、$a=A_4/B,c=A_5/B$ とし、
必要なら $t$ の一次式による移動で $A_5$ が零にならないようにする。
$a,c$ および $B,A_6$ はそれぞれ互いに素である。三重印を一つ含む方向は三係数全ての零点である。
局所座標をその印に移せば二次の $t$ 係数は位数一以上、一次係数は二以上、定数項は三以上になる。
さらに同方向の相異なる三重印数 $k\le3$ に対し、元の三係数はいずれも位数 $k$ 以上である。
方向変数での $j$ 次Taylor項は $t$ の二次以下であり、各印で位数 $3-j$ 以上の根を持つ。
$j<k$ のとき $k(3-j)>2$ だからである。

$B\ne0$ の方向（$B$ の零集合の外）では $I_z(a,c)\ge2k$ である。
$k=1$ では $t$ を移して一次係数の位数二を使い、$k\ge2$ では $k^2\ge2k$ を使う。
$B=0$ の方向で $r=\operatorname{ord}_z B$ とする。
$t$ の移動で変わる定数項は $B$ を法として $A_6$ と等しいので、$I_z(B,A_6)\ge3r$。
$k>r$ なら $I_z(a,c)\ge(k-r)^2$ である。従って $k$ に対して
$I_z(B,A_6)/3+I_z(a,c)/2$ を充てると、不足は $k-r=1$ の場合の高々 $1/2$ だけである。
このような方向は $B\cdot A_6=6b$ より高々 $2b$ 個である。
これで $p$ 以外の印数は
\[
                       3b+(4-b)(5-b)/2
\]
以下となる。$0\le b\le3$ なら十以下である。$b=4$ なら $a$ は単元で常に $k\le r$ となり、
$B\cdot A_6$ のみから八以下である。従って $p$ を印に含めても十三には届かない。
\end{proof}
\begin{proof}[三つの極条件の局所交数]
今やどの三重印も重複度がちょうど三である。その一つ $p$ を座標点にし、
$G=t^3 A_3(z)+t^2 A_4(z)+t A_5(z)+A_6(z)$ と書く。
三曲面 $G,\partial_tG,\partial_t^2G$ が固有に交わると仮定すれば全長は $6\cdot5\cdot4=120$ である。
$p$ の局所座標 $t=1$ では三式の $A_3,A_4,A_5$ の係数行列は可逆である。
行基本操作により位数三、四、五以上の三生成元が得られ、局所交数は少なくとも六十である。
他の三重印では微分で位数は一ずつしか落ちず、交数は少なくとも六である。
十二印を加えれば $60+12\cdot6>120$。従って共通零集合は正次元である。
$G$ は整で一階極は次数五の非零式なので、共通零集合の正次元成分は曲線である。
\end{proof}
